% Self-contained source. No external bibliography or figures are required.
% Compile twice with: pdflatex third_scheme_note.tex
\documentclass[11pt]{article}
\usepackage[T1]{fontenc}
\usepackage[margin=24mm]{geometry}
\usepackage{amsmath,amssymb,amsthm}
\usepackage[hidelinks]{hyperref}
\allowdisplaybreaks
\setlength{\parindent}{0pt}
\setlength{\parskip}{0.55em}
\setlength{\emergencystretch}{2em}
\newcommand{\E}{\mathbb E}
\newcommand{\R}{\mathbb R}
\newcommand{\HH}{\mathcal H}
\newcommand{\PP}{\mathcal P}
\newcommand{\Id}{\mathrm{Id}}
\newcommand{\spr}{\operatorname{spr}}
\newcommand{\norm}[1]{\left\lVert #1\right\rVert}
\newtheorem{proposition}{Proposition}
\title{Proposed additions for the third neural scheme}
\author{}
\date{}
\begin{document}
\maketitle
\vspace{-2em}

\emph{Placement: in Section~3.2.2 of~\cite{GRS}, after the introductory
sentence describing the epoch structure and before Algorithm~2. The two
paragraphs below explain the method's relationship to existing approaches.}

Let $\mathcal U_\zeta=(u_\zeta,z_\zeta)$ denote the network approximation
of $(u,\bar u)$, and let $\Phi$ be the nonlinear Feynman--Kac operator
introduced in Section~2.3. The third scheme seeks to satisfy the integral
model constraint $R_\zeta:=\mathcal U_\zeta-\Phi(\mathcal U_\zeta)=0$.
This relates the method to physics-informed neural networks
(PINNs)~\cite{RPK}, which impose a model equation through a residual loss.
Here the residual comes from a probabilistic integral representation;
strong-form PINNs typically use the differential residual
$\mathcal L u_\zeta+f(x,u_\zeta,\nabla u_\zeta\sigma)$.
The Malliavin representation supplies the gradient component through
an expectation, so evaluating the model constraint requires no spatial
derivatives of the network. A close precedent is the derivative-free
PDE method of Han, Nica, and Stinchcombe~\cite{HNS}, which also combines
probabilistic representations with neural regression. The present scheme
uses the coupled representations of the solution and its gradient for
an infinite-horizon BSDE.

The training procedure also connects to fitted value iteration~\cite{MS}
and target-network methods~\cite{Mnih}. At epoch $i$, starting points
$x_k$ are sampled and targets
$Y_k^i=\widehat\Phi_M(\mathcal U_{\bar\zeta_i})(x_k)$ are computed from
the parameters $\bar\zeta_i$ at the start of the epoch. These targets
are held fixed during $\PP$ optimizer updates. With
$\HH=L^2(\mu_0;\R^{d'+d'd})$, the empirical squared loss and its
population counterpart are
\begin{align}
 \widehat Q_i(\zeta)
  &=\frac{1}{2M_x}\sum_{k=1}^{M_x}
           |\mathcal U_\zeta(x_k)-Y_k^i|^2,\notag\\
 Q(\zeta;\vartheta)
  &=\tfrac12\norm{\mathcal U_\zeta-\Phi(\mathcal U_\vartheta)}_\HH^2.
                                                        \label{eq:third-loss}
\end{align}
Thus the method performs regression towards a periodically refreshed
model-based target. This connection also specifies which derivative is
used: writing $J_\zeta=D_\zeta\mathcal U_\zeta$ and
$T_\zeta=D_\zeta[\Phi(\mathcal U_\zeta)]$, one obtains, where differentiable,
\begin{equation}
 \left.\nabla_1 Q(\zeta;\vartheta)\right|_{\vartheta=\zeta}
       =J_\zeta^*R_\zeta,
 \qquad
 \nabla_\zeta\!\left(\tfrac12\norm{R_\zeta}_\HH^2\right)
       =(J_\zeta-T_\zeta)^*R_\zeta.                       \label{eq:third-gradient}
\end{equation}
The stars denote Hilbert adjoints. The frozen-target update uses the
first derivative, so its stability must be studied as a regression
iteration with changing targets.

\emph{Placement: immediately after Algorithm~2 and before Section~4,
replacing the sentence that defers the convergence study entirely to
future work. The proposition states a conditional result for the
population, fixed-preconditioner version of the algorithm.}

Consider an epoch consisting of $\PP$ population gradient updates with
constant learning rate $\eta>0$ and a fixed symmetric positive definite
preconditioner $D$:
\begin{equation}
 \zeta_{i,0}=\bar\zeta_i,\qquad
 \zeta_{i,\ell+1}=\zeta_{i,\ell}
       -\eta D\nabla_1 Q(\zeta_{i,\ell};\bar\zeta_i),
 \qquad \bar\zeta_{i+1}=\zeta_{i,\PP}=: \Psi_\PP(\bar\zeta_i).
                                                        \label{eq:third-epoch}
\end{equation}
Ordinary gradient descent corresponds to $D=\Id$.

\begin{proposition}[Conditional local convergence]\label{prop:third-local}
Assume $\zeta\mapsto\mathcal U_\zeta$ and
$\zeta\mapsto\Phi(\mathcal U_\zeta)$ are $C^2$ as maps into $\HH$ near
$\zeta_*$, where $R_{\zeta_*}=0$. In locally identifiable coordinates,
assume $G=J_*^*J_*\succ0$, and define
\begin{equation}
 C=J_*^*T_*,\qquad B=G^{-1}C,\qquad A=\Id-\eta DG.
                                                        \label{eq:third-matrices}
\end{equation}
Then the derivative of one complete epoch is
\begin{equation}
 S_\PP:=D\Psi_\PP(\zeta_*)
       =A^\PP+\eta\sum_{j=0}^{\PP-1}A^jDC
       =A^\PP+(\Id-A^\PP)B.                             \label{eq:third-jacobian}
\end{equation}
If $\spr(S_\PP)<1$, there exist $W\succ0$, $r>0$, and $q\in(0,1)$ such
that every initialization satisfying
$\norm{\bar\zeta_0-\zeta_*}_W\leq r$ converges according to
\begin{equation}
 \norm{\bar\zeta_i-\zeta_*}_W
       \leq q^i\norm{\bar\zeta_0-\zeta_*}_W,
 \qquad \norm{h}_W=(h^\top Wh)^{1/2}.                    \label{eq:third-rate}
\end{equation}
Here $\spr$ denotes spectral radius.
\end{proposition}
\begin{proof}
At $(\zeta_*,\zeta_*)$, the zero residual gives
$\partial_{11}Q=G$ and $\partial_{12}Q=-C$. An initial perturbation $h$
therefore satisfies $h_0=h$ and
$h_{\ell+1}=Ah_\ell+\eta DC h$. Iteration gives the sum
in~\eqref{eq:third-jacobian}; using $\eta DG=\Id-A$ gives its last
expression. If $\spr(S_\PP)<1$, the discrete Lyapunov equation
$S_\PP^\top W S_\PP-W=-\Id$ admits $W\succ0$.
In the induced norm, $\norm{S_\PP}_W<1$. Continuity of the derivative
makes $\Psi_\PP$ a $q$-contraction on a sufficiently small convex ball
about its fixed point $\zeta_*$. The ball is invariant, and the
contraction theorem gives~\eqref{eq:third-rate}; see also~\cite{Kelley}.
\end{proof}

A neighborhood can also be certified around an approximate learned
candidate $\widehat\zeta$. Fix $W\succ0$ and a closed ball
$\mathcal B_r=\{\zeta:\norm{\zeta-\widehat\zeta}_W\leq r\}$ in the
$C^1$ domain of the epoch map $\Psi=\Psi_\PP$. Suppose
\begin{align}
 \epsilon&=\norm{\Psi(\widehat\zeta)-\widehat\zeta}_W,
 &s&=\norm{D\Psi(\widehat\zeta)}_W,\notag\\
 \norm{D\Psi(\zeta)-D\Psi(\widehat\zeta)}_W
      &\leq L\norm{\zeta-\widehat\zeta}_W\quad(\zeta\in\mathcal B_r),
 &q&=s+Lr<1,\notag\\
 \epsilon+qr&\leq r.                                    \label{eq:third-ball}
\end{align}
These bounds make $\mathcal B_r$ invariant and $\Psi$ a contraction on
it. Every initialization in this ball therefore converges to its unique
fixed point $\zeta^\dagger$, with
$\norm{\zeta^\dagger-\widehat\zeta}_W\leq\epsilon/(1-q)$.
The bounds must hold on the whole ball; the spectrum at the candidate
alone does not certify a region of convergence. Identifying this training
fixed point with a BSDE solution additionally requires $R_{\zeta^\dagger}=0$
and admissibility in the solution class.

This local criterion differs from the global requirement used for the
first two schemes,
\begin{equation}
 \norm{\Phi(w)-\Phi(\widetilde w)}_{\mathcal X}
       \leq\kappa_p\norm{w-\widetilde w}_{\mathcal X}
       \quad\text{for all admissible }w,\widetilde w,
 \qquad \kappa_p<1.                                    \label{eq:third-global}
\end{equation}
The condition $\spr(S_\PP)<1$ only constrains the training dynamics near
one root, and depends on the learning rate, the epoch length, and the
local neural representation. If the global $p=2$ estimate holds in the
same $\HH$ as the loss, then
$\norm{B}_G\leq\kappa_2$, where
$\norm{h}_G^2=h^\top Gh$; this comparison requires the measure-stability
assumptions of that estimate. When $\spr(A)<1$, letting $\PP\to\infty$
in~\eqref{eq:third-jacobian} gives $B$, the derivative of a locally fully
fitted projected Picard step. Finite training blocks change this matrix
and can therefore change its stability.

For the curvature preconditioner $D=G^{-1}$ and $0<\eta<1$, the
linearization takes the relaxed Picard form
\begin{equation}
 S_\PP=(1-\alpha)\Id+\alpha B,
 \qquad \alpha=1-(1-\eta)^\PP.                           \label{eq:third-relaxation}
\end{equation}
A mode $\lambda=1$ is neutral. For
$\lambda\in\operatorname{spec}(B)\setminus\{1\}$, stability is equivalent to
\begin{equation}
 |1-\alpha+\alpha\lambda|<1
 \quad\Longleftrightarrow\quad
 \operatorname{Re}\lambda<1\ \text{ and }
 0<\alpha<\frac{2(1-\operatorname{Re}\lambda)}{|1-\lambda|^2}.
                                                        \label{eq:third-region}
\end{equation}
Consequently finite fitting can stabilize modes with $|\lambda|>1$
whose real part is below $1$. This supplies a local mechanism for
convergence beyond the Picard stability region. The general fixed
preconditioner is governed by~\eqref{eq:third-jacobian}. Failure of a
global upper bound $\kappa_p<1$ by itself establishes neither Picard
instability nor stability of the third scheme.

The result concerns the population squared-loss iteration
\eqref{eq:third-epoch}. Applying it to the implementation's stochastic
Adam updates, RMS objective, and evolving batch-normalization state
requires the corresponding state and sampling analysis. Local smoothness
and identifiability must also be justified for the network being used.
These qualifications specify the scope of the proposition; the existing
numerical results alone do not verify its hypotheses.

\begin{thebibliography}{9}
\bibitem{GRS}
E.~Gobet, A.~Richou, and C.~Shardul.
\emph{Numerical approximation of Markovian BSDEs in infinite horizon
and elliptic PDEs}. arXiv:2604.24392, 2026.
\url{https://arxiv.org/abs/2604.24392}.

\bibitem{RPK}
M.~Raissi, P.~Perdikaris, and G.~E.~Karniadakis.
Physics-informed neural networks: A deep learning framework for solving
forward and inverse problems involving nonlinear partial differential
equations. \emph{Journal of Computational Physics}, 378:686--707, 2019.
\url{https://doi.org/10.1016/j.jcp.2018.10.045}.

\bibitem{HNS}
J.~Han, M.~Nica, and A.~R.~Stinchcombe.
\emph{A derivative-free method for solving elliptic partial
differential equations with deep neural networks}.
arXiv:2001.06145, 2020.
\url{https://arxiv.org/abs/2001.06145}.

\bibitem{MS}
R.~Munos and C.~Szepesv\'ari.
Finite-time bounds for fitted value iteration.
\emph{Journal of Machine Learning Research}, 9:815--857, 2008.
\url{https://jmlr.org/papers/v9/munos08a.html}.

\bibitem{Mnih}
V.~Mnih et al.
Human-level control through deep reinforcement learning.
\emph{Nature}, 518:529--533, 2015.
\url{https://doi.org/10.1038/nature14236}.

\bibitem{Kelley}
C.~T.~Kelley.
\emph{Iterative Methods for Linear and Nonlinear Equations}.
SIAM, 1995.
\url{https://doi.org/10.1137/1.9781611970944}.
\end{thebibliography}
\end{document}
